\section{What the reconstruction guarantees and what it does not}
A smooth common potential couples derivatives everywhere in its model domain.
Its reciprocal mixed slopes and optional reflection do not depend on sample
fit error. Gauge removes a constant, while measured support and regularization
control additional weak directions. A full-rank design does not ensure good
conditioning, boundary accuracy, positive Hessian, or identifiable sensor bias.

Synthetic truth is representable in this basis and uses one seed, simple noise
and prescribed outliers. Its accuracy results do not establish a universal
winner. RBF alternatives are conceptual, not a numerical benchmark.
The real experiment has unknown magnetic truth and uses terminal currents,
nominal resistance and recorded electrical-speed/voltage channels. Its
spatial holdout reduces local leakage but selects the penalty on the reported
score. Additional independent campaigns and calibrated derivative references
would be needed for a physical accuracy claim.

The model predicts consistent slopes. The residual carries incompatible
measurements, excluded physical state, interpolation/basis error and penalty
bias together. Storage-coordinate qualifications are developed in the
companion consistency study \cite{companion-consistency}; here they are a
gate on interpretation rather than another fitting parameter.
